\chapter{Web Service Tasks}
\label{ch:web-tasks}

\section{Design Goals}

\texttt{TaskType.WEB} accepts a JSON-serializable HTTP/HTTPS request
specification \emph{rather than a shell command}. The workflow driver writes
that specification under \texttt{.dagon/}, stages any referenced files below
\texttt{inputs/}, and invokes the web runner inside the task directory.
Consequently, the HTTP request occurs \emph{where the task executes}: a
\texttt{Slurm}-scheduled web task (\texttt{executor="slurm"}) issues its
request from the compute node, not from the controlling process, without
changing the dataflow model in any way.

\begin{practicebox}
The lesson's stated objective captures the design intent precisely: a web
request is a first-class DAGonStar task --- its dependency is inferred from
\texttt{workflow://}, its input is staged into a private scratch directory,
and its response is a declared workflow output, exactly like any other task
type covered in this book.
\end{practicebox}

\begin{lessonbox}{Lecture 26 --- Uploads, authentication and failures}
\textbf{Objective.} Construct a reproducible workflow in which a local HTTP
service receives a file produced by an upstream task.

\textbf{Course example} (\texttt{examples/26-web-upload/web\_upload.py}, Lecture~26; Lecture~25's \texttt{examples/25-web-get/web\_get.py} covers GET requests).
\lstinputlisting[caption={Lecture 26: \texttt{web\_upload.py}},label={lst:lecture-26}]{../../examples/26-web-upload/web_upload.py}

\textbf{Run} from the repository root:
\begin{lstlisting}[style=bash]
python examples/26-web-upload/web_upload.py
\end{lstlisting}

\textbf{Recorded output} (\texttt{examples/26-web-upload/output/web\_upload.txt}):
\lstinputlisting[style=bash]{../../examples/26-web-upload/output/web_upload.txt}

\textbf{Scientific practice note.} A web request is, in this model, exactly
as auditable as a shell command: its specification, staged inputs, and
declared outputs all live inside the task's scratch directory, alongside the
launcher-based provenance developed in \Cref{ch:lessons-03-04}.
\end{lessonbox}

\begin{figure}[h]
\centering
\begin{tikzpicture}[
  node distance=20mm,
  task/.style={rectangle, rounded corners, draw=dagonblue, very thick,
    fill=blue!4, minimum width=30mm, minimum height=11mm, font=\sffamily\small, align=center},
  arr/.style={-{Stealth[length=2.6mm]}, thick, dagonblue}
]
  \node[task] (prep) {fetch\\ \footnotesize \texttt{outputs/observations.csv}};
  \node[task, right=26mm of prep] (upl) {upload\\ \footnotesize \texttt{TaskType.WEB} (multipart POST)};
  \node[task, right=26mm of upl] (con) {receipt\\ \footnotesize \texttt{outputs/receipt.json}};
  \draw[arr] (prep) -- node[above,font=\sffamily\scriptsize]{\texttt{workflow:///fetch/...}} (upl);
  \draw[arr] (upl) -- node[above,font=\sffamily\scriptsize]{response body} (con);
\end{tikzpicture}
\caption{The upload of Lecture~26: a file producer, an HTTP upload step whose
multipart payload is a staged \texttt{workflow://} reference, and the declared
response output.}
\label{fig:web-task-flow}
\end{figure}

\section{Interpretation and Inspection}

The multipart wrapper

\begin{lstlisting}
"multipart": {
    "dataset": {"file": "workflow:///fetch/outputs/observations.csv",
                "content_type": "text/csv"},
    "description": {"value": "Gulf of Naples, 2 Nov 2026"}}
\end{lstlisting}

explicitly distinguishes a file upload from an ordinary string field.
DAGonStar creates the \texttt{fetch} $\to$ \texttt{upload} dependency from
the reference and copies the file into
\texttt{inputs/Upload26/fetch/outputs/observations.csv} inside the upload
task's folder, mirroring exactly
the staging behavior already established for shell-command tasks
(\Cref{ch:schema}). The response body is written to the declared
\texttt{outputs/receipt.json}; response metadata --- including status code,
content hash, and a \emph{non-secret} authentication description such as
\texttt{\{"source": "env:STATION\_TOKEN", "type": "bearer"\}} --- is written
to \texttt{.dagon/web\_result.json}, and the sanitised request to
\texttt{.dagon/web\_request.json}.

\section{Authentication}

For authenticated endpoints, the specification uses environment-variable
\emph{names}, never literal secrets:

\begin{lstlisting}
{"type": "bearer", "token_env": "SERVICE_TOKEN"}
\end{lstlisting}

Basic and API-key authentication modes follow the same principle (see
\texttt{username\_env}, \texttt{password\_env}, \texttt{value\_env} in
\Cref{ch:configuration}).

\begin{cautionbox}
Never place secret values in a workflow, a generated request file, a task
output, or a tutorial. The environment-variable indirection is not optional
boilerplate; it is the load-bearing mechanism that keeps credentials out of
version-controlled workflow definitions and out of checkpoint or
serialization JSON.
\end{cautionbox}

\section{Specification Surface}

Beyond \texttt{multipart}, the full \texttt{WebTask} specification supports
\texttt{query}, \texttt{headers}, \texttt{json}, \texttt{data}, \texttt{auth},
\texttt{timeout}, retry settings, \texttt{expected\_status} codes, and
\texttt{outputs}. Recursive \texttt{workflow://} values anywhere in the
specification infer dependencies; referenced files are copied below
\texttt{inputs/}. Body, headers, and metadata outputs are written to safe
relative paths below \texttt{outputs/}. \texttt{executor="slurm"} accepts
the same scheduler options, in \texttt{resources}, used by Slurm batch tasks
(\Cref{ch:backends}) --- allowing a web request to be issued from a
scheduled compute node, for instance to publish a forecast product from the
same node that produced it.

\section{Verification}

\begin{lstlisting}[style=bash]
python3 -m unittest tests.test_web -v
\end{lstlisting}

\section{Scientific Application: Web Tasks as a Publication Stage}

In the operational forecasting-chain pattern developed fully in
\Cref{ch:case-study}, the Web task type is most naturally the
\emph{publication} stage of a pipeline: a model's post-processed product is
uploaded to a dissemination API, a dashboard, or a data catalog as the final
DAG node, with the same dependency-graph guarantees --- it cannot run before
post-processing completes, and its outcome is recorded as an ordinary,
inspectable task output --- as every preceding stage.

\section{Idempotency and the Risks of Retrying a Web Task}

A subtlety that the mechanical treatment of Lecture~25 does not surface, but
that any researcher moving beyond the tutorial's local echo-server example
must reason about explicitly, concerns \emph{idempotency} --- whether
repeating a request has the same effect as issuing it once.

An HTTP \texttt{GET} request is, by the protocol's own design convention,
expected to be idempotent and free of side effects: issuing it twice should
leave the world in the same state as issuing it once, which makes automatic
retry behavior (referenced among the specification surface in
\Cref{ch:web-tasks} as ``retry settings'') generally safe to enable. An
HTTP \texttt{POST} request, by contrast, frequently \emph{is not}
idempotent: a publication endpoint that appends a new forecast record to a
database, or a notification endpoint that sends an email, may have a
materially different effect if invoked twice than if invoked once,
regardless of whether the second invocation was a deliberate workflow
rerun or an automatic retry triggered by a transient network failure.

\begin{cautionbox}
A workflow author configuring a \texttt{WEB} task's retry behavior for a
non-idempotent endpoint --- the publication stage of the forecasting-chain
pattern developed in \Cref{ch:synthesis} being the clearest example in this
book --- should verify, against the target service's own documentation,
whether the endpoint itself provides idempotency guarantees (for instance,
via a client-supplied idempotency key included in the request body or
headers) before relying on the engine's retry mechanism. Where no such
guarantee exists, automatic retries risk silently duplicating a published
scientific product, an error considerably harder to detect after the fact
than an outright request failure would be.
\end{cautionbox}

\section{Comparing Web Tasks and LLM Tasks as Sibling Mechanisms}

Having now examined both \texttt{TaskType.LLM} (\Cref{ch:llm-tasks}) and
\texttt{TaskType.WEB} in full, it is worth making explicit the structural
relationship between them, since a first reading might reasonably ask why
two separate task types exist at all, given that an LLM completion request
is, after all, simply a particular kind of HTTP \texttt{POST}.

The answer is that \texttt{TaskType.LLM} can be understood as a
\emph{specialization} of the general HTTP-request capability
\texttt{TaskType.WEB} provides, pre-configured with the specific
conventions --- the \texttt{/chat/completions} path convention, the
\texttt{messages}-based JSON body shape, the \texttt{\{parameter\}} prompt-
templating mechanism, and the \texttt{provider}-section-based credential
lookup of \Cref{sec:config-llm} --- that make authoring an AI-completion
task substantially more convenient than assembling the equivalent raw
\texttt{WEB} task specification by hand would be. \Cref{tab:ai-hpc-triad} in
\Cref{ch:native-tasks} will complete this comparison with the third sibling,
\texttt{TaskType.NATIVE}; the present comparison can already be made
concrete:

\begin{table}[h]
\centering
\caption{\texttt{LLM} as a convenience specialization of \texttt{WEB}.}
\begin{tabular}{@{}p{0.3\linewidth}p{0.32\linewidth}p{0.3\linewidth}@{}}
\toprule
\textbf{Aspect} & \textbf{\texttt{TaskType.WEB}} & \textbf{\texttt{TaskType.LLM}} \\
\midrule
Request shape & Fully general (\texttt{method}, \texttt{url}, \texttt{headers}, \texttt{json}, \texttt{multipart}, \dots) & Fixed to a Chat Completions JSON body with \texttt{messages} \\
Endpoint configuration & Literal \texttt{url} per task & Resolved from \texttt{[llm.<provider>]} section \\
Templating & None built in; values must be assembled by the caller & \texttt{\{parameter\}} substitution from \texttt{params}/\texttt{input\_files} \\
Typical role & Publication, ingestion, arbitrary third-party APIs & AI-assisted interpretation, summarization \\
\bottomrule
\end{tabular}
\end{table}

\begin{practicebox}
A useful diagnostic question for a workflow author unsure which task type to
reach for: if the target endpoint is an OpenAI-compatible Chat Completions
API, \texttt{TaskType.LLM}'s built-in conventions remove enough boilerplate
to be worth using even though the underlying mechanism --- an HTTP request
with a JSON body --- is the same one \texttt{TaskType.WEB} would also
support; for every other HTTP/HTTPS service, \texttt{TaskType.WEB}'s full
generality is the correct, and only, choice.
\end{practicebox}

\section{Designing a Robust Publication Specification}

Building on \Cref{ch:web-tasks}'s observation that Web tasks are naturally
suited to the publication stage of a forecasting chain, it is worth working
through, in more depth than the tutorial's locally verifiable example
requires, what a production-quality publication task specification should
additionally account for.

\begin{enumerate}[leftmargin=1.8em]
  \item \textbf{Explicit \texttt{expected\_status} checking.} As shown in
    Lecture~25's specification (\texttt{"expected\_status": 200}), a
    publication task should never treat an HTTP-level success (a response
    received at all) as equivalent to an application-level success (the
    endpoint accepted and processed the published product); checking the
    expected status code is the mechanism that distinguishes the two.
  \item \textbf{Bounded \texttt{timeout}.} Publication endpoints, like LLM
    providers (\Cref{ch:llm-tasks}), may become slow or unresponsive under
    load; an unbounded wait on a stalled publication step blocks the
    \texttt{publish} task and, transitively, anything downstream of it,
    indefinitely.
  \item \textbf{Declared \texttt{outputs} sufficient for audit.} Retaining
    both \texttt{body} and \texttt{metadata} outputs, as Lecture~25 does,
    ensures that a later audit of a published product can confirm not only
    \emph{what} was published but \emph{when} and with what response the
    publishing endpoint acknowledged it.
  \item \textbf{Authentication via environment-variable indirection.} As
    \Cref{ch:configuration} emphasized, and as this chapter's own
    discussion of multipart uploads illustrated, no literal credential
    should ever appear in the specification itself.
\end{enumerate}

\section{Chapter Summary}

\texttt{TaskType.WEB} generalizes file-based scientific software integration
(\Cref{ch:external-software}) to network services: any HTTP/HTTPS endpoint
becomes a DAG node whose inputs are staged via \texttt{workflow://} and whose
outputs are ordinary, dependency-trackable files. This chapter additionally
examined the idempotency risks of automatic retries against non-idempotent
publication endpoints, situated \texttt{TaskType.LLM} explicitly as a
convenience specialization of the more general \texttt{TaskType.WEB}
mechanism, and worked through a four-point checklist for designing a
production-quality publication task specification. \Cref{ch:native-tasks}
completes the AI--HPC integration triad by bringing native, importable
Python functions into the same model.

\begin{exercise}
Using the four-point checklist above, critique the minimal Lecture~25
specification reproduced earlier in this chapter, identifying which of the
four points it already satisfies and which would need strengthening before
the same specification could be safely used against a real, non-idempotent,
internet-facing publication endpoint rather than the tutorial's local echo
server.
\end{exercise}
